\documentclass[a4paper,11pt]{article}
\usepackage[T2A]{fontenc}
\usepackage[utf8]{inputenc}
\usepackage[russian]{babel}
\usepackage{amsmath,amssymb,mathtools}
\usepackage{helvet}
\renewcommand{\familydefault}{\sfdefault}
\usepackage{geometry}
\geometry{left=22mm,right=22mm,top=22mm,bottom=22mm}
\usepackage{microtype}
\usepackage{icomma}
\usepackage{tikz}
\usetikzlibrary{arrows.meta,positioning,shapes.geometric,calc}
\usepackage{booktabs,tabularx,array}
\usepackage{enumitem}
\usepackage{parskip}
\usepackage{fancyhdr}
\usepackage{setspace}
\usepackage{xcolor}
\definecolor{accent}{HTML}{008080}
\definecolor{darkaccent}{HTML}{005858}
\definecolor{textgray}{HTML}{333333}
\definecolor{softgray}{HTML}{6C7375}
\usepackage{hyperref}
\hypersetup{colorlinks=true,linkcolor=darkaccent}
\setlength{\parindent}{0pt}
\setlength{\parskip}{4pt}
\setstretch{1.17}
\setlist[itemize]{leftmargin=5.5mm,itemsep=3pt,topsep=2pt}
\setlist[enumerate]{leftmargin=7mm,itemsep=6pt,topsep=4pt}
\emergencystretch=2em
\sloppy
\pagestyle{fancy}
\fancyhf{}
\fancyhead[L]{\small\textcolor{softgray}{Логарифмы: формулы и вычисления}}
\fancyhead[R]{\small\textcolor{softgray}{Денис Ермолаев}}
\fancyfoot[L]{\footnotesize\textcolor{softgray}{26 октября 2026 г.}}
\fancyfoot[R]{\footnotesize\textcolor{softgray}{\thepage}}
\renewcommand{\headrulewidth}{0.3pt}
\renewcommand{\footrulewidth}{0pt}
\setlength{\headheight}{14pt}
\newcommand{\idea}[1]{\par\textcolor{darkaccent}{\textbf{Ключевая идея.}} #1\par}
\newcommand{\namedexample}[1]{\par\medskip\textbf{#1}\par\vspace{-2mm}}
\newcommand{\topic}[1]{\par\medskip{\large\bfseries\color{darkaccent}#1}\par\vspace{1mm}}
\newcommand{\mathcheck}{\(\square\)}
\begin{document}
\thispagestyle{empty}
\begin{center}
\vspace*{25mm}
{\large\color{darkaccent}\MakeUppercase{Математика}}\par
\vspace{19mm}
{\Huge\bfseries Чек-лист}\par
\vspace{14mm}
{\LARGE\color{accent} Логарифмы: формулы\\[2mm] и вычисления}\par
\vspace{19mm}
{\large Подготовка к ЕГЭ по профильной математике}\par
\vfill
{\Large Денис Ермолаев}\par
\vspace{11mm}
{\large 26 октября 2026 г.}\par
\vspace{13mm}
{\normalsize Лёвин Артём Александрович}\\[2mm]
{\normalsize эксклюзивно для Дениса Ермолаева}\par
\vspace{23mm}
\begin{tikzpicture}[x=1cm,y=1cm]
\draw[accent,line width=.85pt] (0,0) .. controls (1.0,.65) and (2.2,-.45) .. (3.4,0)
 .. controls (4.7,.58) and (5.7,-.52) .. (7,0)
 .. controls (8.1,.48) and (9.0,-.48) .. (10.4,0);
\fill[darkaccent] (5.0,-.02) circle (1.1pt);
\end{tikzpicture}
\end{center}
\clearpage

\section*{1. Смысл логарифма и область определения}
\textbf{Определение.} Логарифм \(\log_a b\) --- показатель степени, в которую надо возвести основание \(a\), чтобы получить \(b\):
\[
\boxed{\log_a b=c\quad\Longleftrightarrow\quad a^c=b.}
\]
В действительных числах обязательны условия
\[
\boxed{a>0,\quad a\ne1,\quad b>0.}
\]
Отсюда \(\log_a1=0\) и \(\log_a a=1\). Например, \(\log_2 8=3\), так как \(2^3=8\).

\section*{2. Формулы, которые надо знать наизусть}
Во всех формулах ниже \(a,b,c>0\), \(a\ne1\); дополнительные ограничения указаны отдельно.
\renewcommand{\arraystretch}{1.42}
\begin{tabularx}{\linewidth}{@{}>{\raggedright\arraybackslash}p{.31\linewidth}X@{}}
\toprule
\textbf{Приём} & \textbf{Формула}\tabularnewline
\midrule
Степень аргумента & \(\displaystyle \log_a(b^r)=r\log_a b\)\tabularnewline
Степень основания & \(\displaystyle \log_{a^r}b=\frac1r\log_a b\), \(r\ne0\)\tabularnewline
Сумма логарифмов & \(\displaystyle \log_a b+\log_a c=\log_a(bc)\)\tabularnewline
Разность логарифмов & \(\displaystyle \log_a b-\log_a c=\log_a\frac bc\)\tabularnewline
Основное тождество & \(\displaystyle a^{\log_a b}=b\)\tabularnewline
Переход к основанию \(d\) & \(\displaystyle \log_a b=\frac{\log_d b}{\log_d a}\), \(d>0,\ d\ne1\)\tabularnewline
\bottomrule
\end{tabularx}

\vspace{1mm}
Важное следствие последней формулы:
\[
\boxed{\log_a b=\frac{1}{\log_b a}\qquad(b\ne1).}
\]
\idea{Степень аргумента выносится \emph{без переворота}; степень основания --- с \emph{обратным показателем}. Сумму и разность разрешено объединять только при одинаковых основаниях.}

\section*{3. Степени, корни и дроби}
Прежде чем считать логарифм, запишите аргумент и основание через степени одного числа:
\[
\sqrt[n]{a^m}=a^{m/n}\quad(a>0),\qquad
\frac1{a^r}=a^{-r}.
\]
\namedexample{Пример 1. Дробь и корень из занятия}
\[
\log_3\frac{1}{\sqrt[7]{27}}
=\log_3 3^{-3/7}
=-\frac37\log_3 3=\boxed{-\frac37}.
\]
\namedexample{Пример 2. Степени одновременно в основании и аргументе}
\[
\begin{aligned}
\log_{\frac{1}{\sqrt[5]{125}}}\sqrt[7]{625^4}
&=\log_{5^{-3/5}}5^{16/7}\\
&=\frac{16/7}{-3/5}\log_5 5
=\boxed{-\frac{80}{21}}.
\end{aligned}
\]
Здесь \(125=5^3\), \(625=5^4\). Отрицательный показатель основания нельзя потерять.

\section*{4. Как объединить несколько логарифмов}
Сначала приведите логарифмы к общему основанию. Если перед логарифмом есть числовой коэффициент, перенесите его в степень \emph{аргумента}.
\namedexample{Пример 3. Разные основания}
\[
\begin{aligned}
\log_{27}5+\log_{1/3}2
&=\frac13\log_3 5-\log_3 2\\
&=\log_3\sqrt[3]{5}-\log_3 2\\
&=\boxed{\log_3\frac{\sqrt[3]{5}}2}.
\end{aligned}
\]
Обратите внимание: \(2^{-1}=\frac12\), а не \(-2\).

\section*{5. Число превращаем в логарифм}
Число можно записать в логарифмическом виде с удобным основанием:
\[
\boxed{k=\log_a(a^k)}\qquad(a>0,\ a\ne1).
\]
Например, \(1=\log_2 2\), \(2=\log_3 9\). Это особенно полезно, когда число складывается с логарифмом.
\namedexample{Пример 4. Дробь с логарифмами}
\[
\frac{\log_3 18}{2+\log_3 2}
=\frac{\log_3 18}{\log_3 9+\log_3 2}
=\frac{\log_3 18}{\log_3 18}
=\boxed{1}.
\]

\section*{6. Основное логарифмическое тождество}
Проверяйте, совпадает ли основание внешней степени с основанием логарифма внутри показателя:
\[
3^{\log_3 5}=5,\qquad
3^{2\log_3 5}=3^{\log_3 25}=25.
\]
Если к логарифму прибавлено число, примените свойства степеней:
\[
3^{1+\log_3 5}=3\cdot3^{\log_3 5}=\boxed{15}.
\]
Для другого основания сначала преобразуйте логарифм:
\[
3^{\log_9 5}=3^{\frac12\log_3 5}
=3^{\log_3\sqrt5}=\boxed{\sqrt5}.
\]

\section*{7. Обратные логарифмы и квадрат выражения}
\namedexample{Пример 5. Произведение из занятия}
\[
\begin{aligned}
(1-\log_2 12)(1-\log_6 12)
&=\log_2\frac2{12}\cdot\log_6\frac6{12}\\
&=\log_2\frac16\cdot\log_6\frac12\\
&=(-\log_2 6)(-\log_6 2)\\
&=\log_2 6\cdot\frac1{\log_2 6}=\boxed{1}.
\end{aligned}
\]
Квадрат \emph{всего логарифма} означает квадрат полученного числа:
\[
\bigl(\log_{\sqrt7}49\bigr)^2
=\left(\frac2{1/2}\right)^2=\boxed{16}.
\]
Он отличается от \(\log_{\sqrt7}(49^2)\): это другая запись.

\section*{8. Вычисления с переменными}
\namedexample{Пример 6. Значение выражения через заданный логарифм}
Пусть \(a,b>0\), \(a,b\ne1\), и \(\log_b a=\frac17\). Тогда
\[
\begin{aligned}
\log_a(ab^3)
&=\log_a a+\log_a b^3\\
&=1+3\log_a b\\
&=1+\frac3{\log_b a}=1+3\cdot7=\boxed{22}.
\end{aligned}
\]
Если \(\log_a b=5\), то
\[
\log_a\frac a{b^3}=\log_a a-3\log_a b=1-15=\boxed{-14}.
\]

\section*{9. Прикладная задача: сжатие воздуха}
На занятии использована модель работы \(A\) (в джоулях):
\[
\boxed{A=\alpha\nu T\log_2\frac{V_1}{V_2}},
\]
где \(\alpha\) --- постоянная, \(\nu\) --- количество вещества, \(T\) --- температура, \(V_1,V_2\) --- начальный и конечный объёмы. При \(V_1>0\), \(V_2>0\) отношение положительно, поэтому логарифм определён.

\namedexample{Пример 7. Данные с доски занятия}
При \(A=10350\), \(\alpha=5{,}75\), \(\nu=2\), \(T=300\), \(V_1=40\) найдём \(V_2\):
\[
\begin{aligned}
10350&=5{,}75\cdot2\cdot300\log_2\frac{40}{V_2},\\
\log_2\frac{40}{V_2}&=\frac{10350}{3450}=3,\\
\frac{40}{V_2}&=2^3=8,
\qquad V_2=\boxed{5\ \text{л}}.
\end{aligned}
\]
Проверка: \(V_2<V_1\), что соответствует сжатию.

\section*{10. Типичные ошибки и самопроверка}
\begin{itemize}
\item \(\displaystyle 2^{-1}=\frac12\). Отрицательная степень даёт обратное число.
\item \(\displaystyle \log_{a^r}b=\frac1r\log_a b\), а \(\log_a(b^r)=r\log_a b\). Не путайте роли основания и аргумента.
\item \(\log_a b+\log_c d\) объединяются только после приведения к \emph{одному основанию}.
\item \(\bigl(\log_a b\bigr)^2\) --- квадрат значения логарифма; \(\log_a(b^2)\) --- логарифм квадрата аргумента.
\item При разложении на сумму или разность сначала проверьте, что каждый аргумент положителен.
\end{itemize}

\idea{Для длинного выражения полезен порядок: \textbf{ОДЗ} $\to$ \textbf{степенной вид} $\to$ \textbf{общее основание} $\to$ \textbf{объединение} $\to$ \textbf{вычисление}.}

\section*{Мини-чек-лист перед самостоятельной работой}
Отметьте навыки, которые можете применить без подсказки.
\begin{itemize}[label=$\square$,leftmargin=8mm,itemsep=5pt]
\item Могу объяснить смысл \(\log_a b\) и записать ограничения на основание и аргумент.
\item Различаю степень аргумента, степень основания и степень самого логарифма.
\item Умею привести логарифмы к одному основанию и объединить сумму или разность.
\item Знаю, как представить число логарифмом и применить \(a^{\log_a b}=b\).
\item Узнаю взаимно обратные логарифмы и не теряю знак при отрицательных степенях.
\item Могу найти неизвестный объём по формуле, содержащей логарифм отношения объёмов.
\end{itemize}

\clearpage
\section*{Алгоритм: как упрощать выражения с логарифмами}
\vspace{6mm}
\begin{center}
\tikzset{
 flowstep/.style={rectangle,draw=darkaccent,very thin,rounded corners=1.6pt,text width=10.2cm,minimum height=12mm,align=center,inner sep=5pt,font=\small},
 decision/.style={diamond,draw=darkaccent,aspect=3.1,text width=4.6cm,align=center,inner sep=1.7pt,font=\small},
 terminal/.style={ellipse,draw=darkaccent,thin,minimum width=7.1cm,minimum height=10mm,align=center,font=\small},
 arrow/.style={-Latex,draw=darkaccent,line width=.7pt}
}
\begin{tikzpicture}[node distance=10mm]
\node[terminal] (start) {Дано выражение с логарифмами};
\node[flowstep,below=of start] (odz) {Проверьте ОДЗ: основания положительны и не равны 1; аргументы положительны};
\node[flowstep,below=of odz] (powers) {Представьте составные числа, корни и дроби в степенном виде};
\node[decision,below=10mm of powers] (different) {Есть логарифмы с разными основаниями?};
\node[flowstep,below=16mm of different] (common) {Приведите их к общему основанию по формуле перехода};
\node[flowstep,below=of common] (coeff) {Перенесите коэффициенты перед логарифмами в степени аргументов};
\node[flowstep,below=of coeff] (combine) {Объедините суммы и разности по формулам произведения и частного};
\node[flowstep,below=of combine] (finish) {Примените обратные логарифмы, тождество $a^{\log_a b}=b$ и определения};
\node[terminal,below=of finish] (end) {Вычислите результат и проверьте ограничения};
\draw[arrow] (start)--(odz);
\draw[arrow] (odz)--(powers);
\draw[arrow] (powers)--(different);
\draw[arrow] (different)--node[right,font=\footnotesize] {да}(common);
\coordinate (detour) at ($(different.east)+(2.8,0)$);
\draw[arrow] (different.east)--node[above,font=\footnotesize] {нет}(detour) |- (coeff.east);
\draw[arrow] (common)--(coeff);
\draw[arrow] (coeff)--(combine);
\draw[arrow] (combine)--(finish);
\draw[arrow] (finish)--(end);
\end{tikzpicture}
\end{center}
\clearpage

\section*{Тренировочный блок --- Путь А}
\textcolor{softgray}{10 задач по возрастанию сложности. Решения записывайте последовательно; ответы расположены на следующей странице.}
\topic{Базовый уровень}
\begin{enumerate}[label=\textbf{\arabic*.}]
\item Вычислите: \(\displaystyle \log_4 64\).
\item Вычислите: \(\displaystyle \log_3\frac1{\sqrt[5]{81}}\).
\item Вычислите: \(\displaystyle \log_{\frac1{\sqrt[3]{16}}}\sqrt[5]{64}\).
\end{enumerate}
\topic{Уровень выше среднего}
\begin{enumerate}[label=\textbf{\arabic*.},start=4]
\item Представьте одним логарифмом: \(\displaystyle \log_2 7+\log_8 9\).
\item Вычислите: \(\displaystyle 5^{2+\log_5 3}\).
\item Вычислите: \(\displaystyle (1-\log_3 15)(1-\log_5 15)\).
\item Вычислите: \(\displaystyle \bigl(\log_{\sqrt{11}}121\bigr)^2\).
\item Вычислите: \(\displaystyle \frac{\log_4 80}{2+\log_4 5}\).
\end{enumerate}
\topic{Сложный уровень}
\begin{enumerate}[label=\textbf{\arabic*.},start=9]
\item Пусть \(a,b>0\), \(a\ne1\), \(b\ne1\), \(\log_b a=\frac13\). Найдите \(\displaystyle \log_a(ab^4)\).
\item В модели \(\displaystyle A=\alpha\nu T\log_2\frac{V_1}{V_2}\) известно: \(A=4500\) Дж, \(\alpha=3\), \(\nu=2\), \(T=250\) К и \(V_1=48\) л. Найдите конечный объём \(V_2\) (в литрах).
\end{enumerate}
\clearpage
\section*{Ответы к тренировочному блоку}
\vspace{4mm}
\renewcommand{\arraystretch}{1.65}
\begin{center}
\begin{tabularx}{.90\linewidth}{@{}>{\bfseries}p{17mm}X@{}}
\toprule
№ & Ответ \\
\midrule
1 & \(3\) \\
2 & \(-\dfrac45\) \\
3 & \(-\dfrac9{10}\) \\
4 & \(\displaystyle \log_2\bigl(7\sqrt[3]{9}\bigr)\) \\
5 & \(75\) \\
6 & \(1\) \\
7 & \(16\) \\
8 & \(1\) \\
9 & \(13\) \\
10 & \(6\) л \\
\bottomrule
\end{tabularx}
\end{center}
\vspace{10mm}
\section*{Итоговая памятка}
\begin{tabularx}{\linewidth}{@{}p{.42\linewidth}X@{}}
\toprule
\textbf{Что увидели} & \textbf{Что сделать}\\
\midrule
Корень, дробь, составное число & Перепишите в степенном виде.\\
Степень у аргумента & Вынесите показатель без переворота.\\
Степень у основания & Вынесите обратный показатель.\\
Сумма или разность логарифмов & Приведите к одному основанию, затем объедините.\\
Число рядом с логарифмом & Запишите число как логарифм удобного основания.\\
Логарифмы с переставленными числами & Проверьте формулу взаимной обратности.\\
\bottomrule
\end{tabularx}
\vfill
\begin{center}
\textcolor{softgray}{\small Лёвин Артём Александрович --- эксклюзивно для Дениса Ермолаева}
\end{center}
\end{document}
